\section{Exact saddle expansion and polynomial
perturbations}\label{exact-saddle-expansion-and-polynomial-perturbations}

\subsection{Definitions and statement}\label{definitions-and-statement}

Put

\[
 M(t)=\frac{t(t-1)}2,\qquad
 H_n=e^{-1}\sum_{m\ge0}\frac{(M(m))_n}{m!},\qquad H_0=1,
\]

where \((u)_n=u(u-1)\cdots(u-n+1)\). For each fixed \(n\), this
numerical sum converges absolutely, since its numerator is a polynomial
of degree \(2n\) in \(m\). Because \(M(m)\) is a nonnegative integer,
terms with \(M(m)<n\) vanish. For \(n\ge1\), define

\[
 a_j=\frac{1+\sqrt{1+8j}}2,\quad b_j=1-a_j,
 \quad f_n(t)=\sum_{j=0}^{n-1}\log(M(t)-j)-\log\Gamma(t+1)
 \quad(t>a_{n-1}).
\]

There is a unique maximum \(\tau=\tau_n\) of \(f_n\). Define

\[
 r=W(2n),\quad s=\frac{2n}{r}=e^r,\quad
 \sigma^2=-\frac1{f_n''(\tau)},\quad
 \lambda_j=f_n^{(j)}(\tau)\sigma^j\quad(j\ge3).
\]

For \(\ell\ge0\), let

\[
 C_\ell=\sum_{\substack{k_3,\ldots,k_{2\ell+2}\ge0\\
               \sum_{j=3}^{2\ell+2}(j-2)k_j=2\ell}}
 \left(\sum_j jk_j-1\right)!!
 \prod_{j=3}^{2\ell+2}\frac{(\lambda_j/j!)^{k_j}}{k_j!},
\]

with the empty summand giving \(C_0=1\) and \((-1)!!=1\).

\textbf{Theorem S1.} For every fixed integer \(K\ge0\),

\[
 H_n=e^{f_n(\tau)-1}\sqrt{2\pi}\,\sigma
       \left(\sum_{\ell=0}^K C_\ell+O_K(n^{-K-1})\right).
 \tag{S1}
\]

Moreover,

\[
 \tau=s+O(r),\qquad
 \sigma^2=\frac{s}{r+1}\left(1+O(r^2/n)\right),\qquad
 \lambda_j=O_j(n^{1-j/2}),\qquad C_\ell=O_\ell(n^{-\ell}).
 \tag{S2}
\]

In particular the error in (S1) has no unreported positive power of
\(r\).

The first two corrections are

\[
 C_1=\frac{\lambda_4}{8}+\frac{5\lambda_3^2}{24},
\]

\[
 C_2=\frac{\lambda_6}{48}+\frac{7\lambda_3\lambda_5}{48}
 +\frac{35\lambda_4^2}{384}+\frac{35\lambda_3^2\lambda_4}{64}
 +\frac{385\lambda_3^4}{1152}.
 \tag{S3}
\]

These coefficients use derivatives at the exact saddle; replacing the
saddle by \(s\) without expanding all accompanying factors changes the
coefficients.

\subsection{Strict concavity, localization, and derivative
scale}\label{strict-concavity-localization-and-derivative-scale}

The factorization

\[
 M(t)-j=\tfrac12(t-a_j)(t-b_j)
\]

gives

\[
 f_n'(t)=\sum_{j<n}\left(\frac1{t-a_j}+\frac1{t-b_j}\right)-\psi(t+1),
\]

\[
 f_n''(t)=-\sum_{j<n}\left(\frac1{(t-a_j)^2}+\frac1{(t-b_j)^2}\right)
 -\psi'(t+1)<0. \tag{S4}
\]

The trigamma function is positive on this interval. Also
\(f_n'(t)\to+\infty\) at the lower endpoint and \(f_n'(t)\to-\infty\) at
infinity. This proves existence and uniqueness of \(\tau\), with no
assumption about the local approximation.

Uniformly for \(t\asymp s\), Taylor expansion of the rational factors
and the usual positive-real digamma estimates yield

\[
 f_n'(t)=\frac{2n}{t}-\log t
     +O\left(\frac{n}{t^2}+\frac{n^2}{t^3}+\frac1t\right), \tag{S5}
\]

\[
 f_n''(t)=-\frac{2n}{t^2}-\frac1t
     +O\left(\frac{n}{t^3}+\frac{n^2}{t^4}+\frac1{t^2}\right). \tag{S6}
\]

Here \(\max_j|a_j|+|b_j|=O(\sqrt n)=o(s)\). At \(s/2\) and \(2s\), the
leading part of (S5) has opposite signs and size comparable to \(r\),
while the error is \(o(r)\). Hence \(\tau\in(s/2,2s)\). At \(s\), the
leading part vanishes and

\[
 f_n'(s)=O(r^3/n),\qquad |f_n''(t)|\asymp n/s^2\asymp r^2/n
\]

on the interval joining \(s\) to \(\tau\). The mean value theorem gives
\(\tau-s=O(r)\). Substituting this into (S6) gives the variance estimate
in (S2).

More generally, for each fixed \(q\ge2\),

\[
 f_n^{(q)}(t)=(-1)^{q-1}(q-1)!
 \sum_{j<n}\left((t-a_j)^{-q}+(t-b_j)^{-q}\right)
 -\psi^{(q-1)}(t+1). \tag{S7}
\]

The polygamma series bounds its last term by \(O_q(t^{1-q})\). Thus,
uniformly on \(I_n=[\tau/2,3\tau/2]\),

\[
 |f_n^{(q)}(t)|\le C_q n/s^q,\qquad
 -f_n''(t)\asymp n/s^2\asymp\sigma^{-2}. \tag{S8}
\]

Since \(s(r+1)\asymp n\), multiplying (S8) by \(\sigma^q\) gives

\[
 \sup_{t\in I_n}|f_n^{(q)}(t)|\sigma^q
 \le C_q n^{1-q/2}. \tag{S9}
\]

This is the uniform derivative estimate responsible for the absence of
logarithmic leakage. It also proves \(C_\ell=O(n^{-\ell})\): each
monomial has order \(n^{-\frac12\sum(j-2)k_j}=n^{-\ell}\).

\subsection{Global tails and a smooth lattice
cutoff}\label{global-tails-and-a-smooth-lattice-cutoff}

Integrating (S8) twice from the saddle gives, for \(t\in I_n\),

\[
 f_n(t)-f_n(\tau)\le-c(t-\tau)^2/\sigma^2. \tag{S10}
\]

In particular, at \(3\tau/4\) and \(5\tau/4\), the loss is at least
\(c n\), because \(\tau^2/\sigma^2\asymp n\). To the left, concavity and
monotonicity bound the entire remaining interval by the left boundary
value. To the right of \(3\tau/2\), concavity gives the tangent bound

\[
 f_n(t)\le f_n(3\tau/2)-c(n/\tau)(t-3\tau/2),\quad t\ge3\tau/2. \tag{S11}
\]

Here (S8), integrated once, shows \(f_n'(3\tau/2)\le-cn/\tau\).
Equations (S10) and (S11) show that both the integral and the integer
sum outside \([3\tau/4,5\tau/4]\) are

\[
 O\left((1+\tau)e^{f_n(\tau)-cn}\right). \tag{S12}
\]

For the sum in the far right tail, sum the exponential tangent majorant
as a geometric series; the far left contains \(O(\tau)\) integer points.
This supplies a global argument, rather than extending a local quadratic
estimate beyond its domain.

Choose a fixed smooth cutoff \(\chi\), equal to one on \([3/4,5/4]\),
supported in \((1/2,3/2)\), with \(0\le\chi\le1\), and put

\[
 g_n(t)=\chi(t/\tau)e^{f_n(t)},
\]

extended by zero to \(\mathbb R\). For every fixed integer \(L\ge0\),

\[
 \|g_n^{(L)}\|_1\le C_L e^{f_n(\tau)}\sigma^{1-L}. \tag{S13}
\]

To see this, write \(x=(t-\tau)/\sigma\). On the support, (S8) and
\(f_n'(\tau)=0\) imply

\[
 |f_n'(t)|\le C\sigma^{-1}|x|,\qquad
 |f_n^{(q)}(t)|\le C_q\sigma^{-q}\quad(q\ge2).
\]

The derivative formula for \(e^{f_n}\) therefore bounds its \(L\)-th
derivative by \(C_L\sigma^{-L}(1+|x|)^L e^{f_n(t)}\). The cutoff
derivatives cost only \(O(\tau^{-q})=O(\sigma^{-q})\). Integrating the
resulting polynomial times the Gaussian majorant (S10) proves (S13).

For even \(L\ge2\), the compact-support Euler--Maclaurin remainder gives

\[
 \left|\sum_{m\in\mathbb Z}g_n(m)-\int_{\mathbb R}g_n(t)\,dt\right|
 \le\frac{2\zeta(L)}{(2\pi)^L}\|g_n^{(L)}\|_1. \tag{S14}
\]

Dividing by \(e^{f_n(\tau)}\sigma\) leaves \(O_L(\sigma^{-L})\). For the
expansion through \(n^{-K}\), take

\[
 L=2K+4.
\]

Then

\[
 \sigma^{-L}=O\left(r^{2K+4}n^{-K-2}\right)=O_K(n^{-K-1}), \tag{S15}
\]

since every fixed power of \(\log n\) is \(O(n)\). Taking only
\(L=2K+2\) would leave \(r^{2K+2}n^{-K-1}\), which does \textbf{not}
justify the claimed remainder. Together with (S12), (S14)--(S15)
rigorously replace the lattice sum by the continuous integral at the
required accuracy.

\subsection{Integrated Taylor expansion to every fixed
order}\label{integrated-taylor-expansion-to-every-fixed-order}

Set \(\eta=n^{-1/2}\) and \(Q=2K+1\). Taylor's theorem and (S9), through
derivative order \(2K+4\), give

\[
 f_n(\tau+\sigma x)-f_n(\tau)
 =-x^2/2+h(\eta,x)+R_n(x),
\]

\[
 h(u,x)=\sum_{j=3}^{2K+3}a_j u^{j-2}x^j/j!,\qquad
 a_j=\lambda_j/\eta^{j-2}=O_j(1),
\]

\[
 |R_n(x)|\le C_K\eta^{2K+2}|x|^{2K+4}. \tag{S16}
\]

These estimates hold throughout \(|x|\le\delta/\eta\), with a
sufficiently small fixed \(\delta>0\) chosen so the associated
\(t\)-interval lies inside the region where the cutoff is one. Reducing
\(\delta\) if necessary, uniformly for \(0\le u\le\eta\),

\[
 |h(u,x)|\le x^2/8,\qquad |R_n(x)|\le x^2/8. \tag{S17}
\]

Indeed each term of \(h\) is bounded by a constant times
\(x^2(\eta|x|)^{j-2}\), and the remainder by a constant times
\(x^2(\eta|x|)^{2K+2}\).

For fixed \(x\), expand \(e^{h(u,x)}\) at \(u=0\) through degree \(Q\).
The coefficients are precisely the monomials indexed by
\(\sum(j-2)k_j\le Q\). The derivative of order \(Q+1\) with respect to
\(u\) is, by repeated differentiation, \(e^{h(u,x)}\) times a polynomial
in derivatives of \(h\). Because all \(a_j\) are bounded and
\(0\le u\le\eta\le1\), this is bounded by

\[
 C_K(1+|x|)^{D_K}e^{x^2/8}
\]

for a finite fixed \(D_K\). The Taylor remainder is therefore bounded by
that expression times \(\eta^{Q+1}\). Also, (S16)--(S17) and the mean
value bound for the exponential imply

\[
 |e^{h(\eta,x)+R_n(x)}-e^{h(\eta,x)}|
 \le C_K\eta^{2K+2}|x|^{2K+4}e^{x^2/4}.
\]

Multiplication by \(e^{-x^2/2}\) gives an integrable polynomial times a
fixed Gaussian. Hence the \textbf{integrated} total remainder is
\(O_K(\eta^{2K+2})=O_K(n^{-K-1})\), without enlarging the error by the
length of the central interval.

Outside \(|x|\le\delta/\eta\), (S10) bounds the cutoff integral by
\(e^{-c_K n}\) times its leading scale. We may likewise extend each of
the finitely many polynomial-Gaussian integrals to all of \(\mathbb R\),
with an exponentially small error. A monomial indexed by \(k_j\) has
ordinary degree \(D=\sum jk_j\) and weighted degree \(w=\sum(j-2)k_j\).
Their parities agree. All odd weighted degrees therefore integrate to
zero. For \(w=2\ell\), the normalized Gaussian moment is \((D-1)!!\),
giving precisely \(C_\ell\).

Thus

\[
 \int g_n(t)\,dt=e^{f_n(\tau)}\sqrt{2\pi}\sigma
 \left(\sum_{\ell=0}^K C_\ell+O_K(n^{-K-1})\right).
\]

Combining this with (S12)--(S15) and the defining factor \(e^{-1}\)
proves Theorem S1.

\subsection{Growth and an adjacent-ratio bound from
positivity}\label{growth-and-an-adjacent-ratio-bound-from-positivity}

\textbf{Proposition S2.}

\[
 \log H_m=2m\log m-2m\log\log m+O(m), \tag{S18}
\]

and, for all sufficiently large \(m\),

\[
 \frac{H_{m-1}}{H_m}\le C\left(\frac{\log m}{m}\right)^2. \tag{S19}
\]

For (S18), set \(r=W(2m)\), \(s=2m/r\), and evaluate the exponent at
\(s\). Since \(m/s^2=o(1)\),

\[
 \sum_{j<m}\log(M(s)-j)=2m\log s-m\log2+O(r^2).
\]

Stirling's estimate, \(s\log s=2m\), and \(\tau-s=O(r)\) give

\[
 f_m(\tau)=2m\log s-m\log2-2m+s+O(r^2+\log m).
\]

Theorem S1 contributes only \(O(\log m)\) to the logarithm. Finally
\(\log s=\log m-\log\log m+O(1)\), giving (S18).

For the ratio, define the positive probability law

\[
 \mathbb P_k(X=m)=\frac{(M(m))_k}{eH_k m!}
\]

on the integers with \(M(m)\ge k\). The exact falling-factorial identity
gives

\[
 \frac{H_{k+1}}{H_k}=\mathbb E_k[M(X)-k]. \tag{S20}
\]

The global tail estimates and the leading case of Theorem S1 show
\(\mathbb P_k(\tau_k/2\le X\le3\tau_k/2)=1-O(e^{-c k})\), after harmless
polynomial factors are absorbed into the exponential. On this interval,

\[
 M(X)-k\ge c\tau_k^2\ge c'k^2/(\log k)^2
\]

for large \(k\). Hence (S20) is at least this size, proving (S19) with
\(m=k+1\). This argument does not subtract two \(O(m)\) logarithmic
errors, which would not establish an adjacent-ratio estimate.

For a more precise conclusion, fix any \(0<\varepsilon<1/2\). Equation
(S10) makes the probability outside \(|X-\tau_k|\le\varepsilon\tau_k\),
but inside \([\tau_k/2,3\tau_k/2]\), exponentially small. Its
contribution to \(\mathbb E_k X^2\) is also exponentially small relative
to \(\tau_k^2\). For \(X>3\tau_k/2\), multiply the exponential tangent
majorant (S11) by \(X^2\) and sum the resulting geometric series and its
first two derivatives; this bounds the second-moment tail by a
polynomial in \(\tau_k\) times \(e^{-ck}\), after normalization. The
left tail is bounded similarly since \(X^2\le\tau_k^2\) there. Thus the
central interval gives
\((1-\varepsilon)^2\le\liminf\mathbb E_kX^2/\tau_k^2\le\limsup\mathbb E_kX^2/\tau_k^2\le(1+\varepsilon)^2\).
Letting \(\varepsilon\downarrow0\), and using \(k=o(\tau_k^2)\), (S20)
gives \(H_{k+1}/H_k\sim\tau_k^2/2\). Consequently

\[
 H_{n-1}/H_n\sim r^2/(2n^2). \tag{S21}
\]

Only the upper bound (S19) is needed below.

A useful explicitly normalized consequence is the following. Define

\[
 S_n=\frac{\exp\{2n\log s-n\log2-2n+s-r^2/4-r/2-1\}}{\sqrt{r+1}}.
\]

Then

\[
 H_n=S_n\left(1+O(r^4/n)\right). \tag{S21a}
\]

Indeed, with \(u_j=1/s+2j/s^2\),

\[
 \sum_{j<n}\log(1-u_j)
 =-\frac ns-\frac{n(n-1)}{s^2}+O\left(\sum_{j<n}u_j^2\right)
 =-r/2-r^2/4+O(r^4/n).
\]

Here \(\max u_j=O(r^2/n)=o(1)\) and \(\sum u_j^2=O(r^4/n)\).
Positive-real Stirling expansion therefore gives

\[
 f_n(s)=2n\log s-n\log2-2n+s-r^2/4-r/2
       -\tfrac12\log(2\pi s)+O(r^4/n).
\]

The estimates \(f_n'(s)=O(r^3/n)\), \(\tau-s=O(r)\), and (S8) imply
\(f_n(\tau)-f_n(s)=O(r^4/n)\). Finally combine the leading case of (S1)
with the variance estimate in (S2). Since \(r^4/n\to0\), exponentiating
the resulting logarithmic error proves (S21a).

\subsection{Arbitrary fixed real polynomial
perturbations}\label{arbitrary-fixed-real-polynomial-perturbations}

Let \(P(x)\) be a fixed real polynomial with \(P(0)=0\), and set

\[
 p_j=[x^j]e^{P(x)},\qquad
 A_P(n)=\sum_{j=0}^n(n)_j p_j H_{n-j}.
\]

This is a finite coefficient identity for formal exponential generating
functions. Neither \(\sum H_n x^n/n!\) nor its product by \(e^{P(x)}\)
is being evaluated at positive \(x\). Equation (S18), together with
Stirling's estimate for \(n!\), shows the former has radius of
convergence zero.

\textbf{Theorem S3.} For each fixed integer \(K\ge0\),

\[
 A_P(n)=\sum_{j=0}^K(n)_j p_jH_{n-j}
       +O_{P,K}\left(H_n(r^2/n)^{K+1}\right). \tag{S22}
\]

If \(P\ne0\) has degree \(d\ge1\), the signed far tail admits the
absolute bound

\[
 \frac1{H_n}\sum_{j>n/2}\left|(n)_j p_jH_{n-j}\right|
 \le\exp\left[-\frac{1+1/d}{2}n\log n+n\log\log n+O_P(n)\right]. \tag{S23}
\]

It also satisfies the weaker bound with leading term \(-n\log n/(2d)\).
For \(P=0\), the tail is identically zero.

\textbf{Proof.} For \(j\le n/2\), multiply (S19) over the indices
\(n-j+1,\ldots,n\):

\[
 \frac{H_{n-j}}{H_n}\le\left(C\frac{\log^2n}{n^2}\right)^j.
\]

Since \((n)_j\le n^j\), these summands are at most \(|p_j|z_n^j\), where
\(z_n=C\log^2n/n\). Write \(P(x)=\sum_{\ell=1}^dc_\ell x^\ell\). The
nonnegative coefficients of \(\exp(\sum |c_\ell|x^\ell)\) majorize
\(|p_j|\). Their series is entire, so, for fixed \(K\),

\[
 \sum_{j\ge K+1}|p_j|z_n^j=O_{P,K}(z_n^{K+1}). \tag{S24}
\]

This bounds the entire part \(K<j\le n/2\).

For the far tail, Cauchy's coefficient inequality on \(|z|=j^{1/d}\)
gives

\[
 |p_j|\le\exp\left[-(j/d)\log j+O_P(j)\right]. \tag{S25}
\]

For completeness, this follows since
\(\max_{|z|=R}|e^{P(z)}|\le\exp(\sum|c_\ell|R^\ell)\), and at
\(R=j^{1/d}\) the last exponent is \(O_P(j)\).

Equation (S18) implies the uniform upper bound

\[
 \log H_k\le2k(\log n-\log\log n)+O(n)\quad(0\le k\le n), \tag{S26}
\]

because \(\log x-\log\log x\) is increasing for \(x>e\), while finitely
many small \(k\) can be absorbed into the constant. Its lower bound at
\(n\), together with \((n)_j\le n^j\), gives, for \(k=n-j\),

\[
 \log\frac{(n)_jH_{n-j}}{H_n}
 \le-j\log n+2j\log\log n+O(n). \tag{S27}
\]

For \(j>n/2\), \(\log j\ge\log n-\log2\), so (S25) and (S27) bound the
logarithm of the absolute normalized summand by

\[
 -(1+1/d)j\log n+2j\log\log n+O_P(n).
\]

The explicit part is decreasing in \(j\) for large \(n\). At \(j=n/2\)
it equals the exponent in (S23); summing at most \(n\) terms only adds
\(O(\log n)\). This proves (S23), which is smaller than every fixed
negative power of \(n\), and (S24) proves (S22), since
\(r\asymp\log n\).

The proof uses an absolute majorant and therefore applies unchanged when
coefficients of \(P\), or the coefficients \(p_j\), have mixed signs. In
particular (S22) with \(K=0\) shows \(A_P(n)/H_n=1+O_P(r^2/n)\), so
\(A_P(n)>0\) for all sufficiently large \(n\), even for signed
perturbations. Consequently its formal exponential generating function
also has radius of convergence zero.

\subsection{A pure power error for the perturbed exact saddle formula}
For $m\ge1$, let
\[
\mathcal H_{m,K}=e^{f_m(\tau_m)-1}\sqrt{2\pi}\,\sigma_m
                   \sum_{\ell=0}^{K}C_{m,\ell}
\]
be the exact-saddle expression truncated at level $K$. For every fixed
$K\ge0$ and fixed $P$, the preceding theorems imply
\[
A_P(n)=\sum_{j=0}^{K+1}(n)_jp_j\mathcal H_{n-j,K}
       +O_{P,K}(H_n n^{-K-1}).
\]
Here $n>K+1$, as is automatic in the limit. Indeed, Theorem S3 with
$K+1$ leaves $O(H_n(r^2/n)^{K+2})=O(H_nn^{-K-1})$, since a fixed
power of $r$ is $o(n)$. On the finitely many retained shifts, Theorem S1
and the adjacent-ratio bound show that the total error of replacing
$H_{n-j}$ by $\mathcal H_{n-j,K}$ is also $O(H_nn^{-K-1})$.
This obtains a pure power error by retaining one extra fixed shift. It
does not turn the same fixed Lambert truncation into a pure power error.
